%% ma: L12-B06-0032
%% lop: 12
%% chuong: 2
%% bai: 6
%% dang: 12.6.13
%% loai: DS
%% muc: [NB, TH, TH, TH]
%% dap_an: "ĐĐĐĐ"
%% nguon: "LaTeX giáo viên gửi 21/10/2026 (M05, Vectơ và hệ trục tọa độ trong không gian, 107 câu), Phần II Câu 4"
%% kien_thuc: "Bài 6 (tích vô hướng, bình phương vô hướng)"
%% trang_thai: cho-duyet
%% ly_do: "Dạng mới đề xuất 12.6.13"
\begin{ex}
Trong không gian, cho hai vectơ $\vec{a}$ và $\vec{b}$ đều có độ dài bằng $1$ và góc giữa chúng là $45^\circ$.
\choiceTF
{\True $\vec{a}\cdot\vec{b}=\dfrac{\sqrt2}2$.}
{\True $|\vec{a}+\vec{b}|^2=2+\sqrt2$.}
{\True $(\vec{a}+3\vec{b})\cdot(2\vec{a}-\vec{b})=-1+\dfrac{5\sqrt2}2$.}
{\True Vectơ $\vec{u}=\vec{a}-\vec{b}$ vuông góc với vectơ $\vec{v}=\vec{a}+\vec{b}$.}
\loigiai{a) $\vec{a}\cdot\vec{b}=1\cdot1\cdot\cos45^\circ=\dfrac{\sqrt2}2$. Đúng.

b) $|\vec{a}+\vec{b}|^2=1+2\cdot\dfrac{\sqrt2}2+1=2+\sqrt2$. Đúng.

c) $(\vec{a}+3\vec{b})\cdot(2\vec{a}-\vec{b})=2|\vec{a}|^2+5\vec{a}\cdot\vec{b}-3|\vec{b}|^2=-1+\dfrac{5\sqrt2}2$. Đúng.

d) $\vec{u}\cdot\vec{v}=|\vec{a}|^2-|\vec{b}|^2=0$ nên $\vec{u}\perp\vec{v}$. Đúng.}
\end{ex}
